\begin{claim}{D.69}{Proved}{A negative quartic coordinate forces a full rotation image}
Use the actual quartic representative, adapted basis, and dihedral
quotient of D.68. For every odd prime $q$ at which $\lambda_q=-1$,
write
\[
 \rho:F\twoheadrightarrow Q=D_{2q},\qquad
 N=\ker\rho,\qquad \beta=\alpha^2,\qquad
 \sigma=s^2,\qquad C=C_F(\sigma),\qquad C_0=C\cap N.
                                                               \tag{D.69a}
\]
The specified rotations form $R=\langle\rho(y)\rangle\cong C_q$.
Then restriction of the actual quotient map gives the exact
sequence and dimension
\[
 1\longrightarrow C_0\longrightarrow C
       \xrightarrow{\ \rho|_C\ }R\longrightarrow1,\qquad
                  \dim_{\F_2}H^1(C_0,\F_2)=q-1.             \tag{D.69b}
\]
The cover defining $N$ is the full completion of an actual
surface of type $(q,2)$.
There are exactly two $N$-conjugacy classes of order-two
sections in $N\rtimes\langle\sigma\rangle$; representatives are
$\langle\sigma\rangle$ and
$\langle x\sigma x^{-1}\rangle$.
Their exact normalizers are
\[
 C_0\times\langle\sigma\rangle,\qquad
 x\bigl(C_0\times\langle\sigma\rangle\bigr)x^{-1}.
                                                               \tag{D.69c}
\]
They are also the two order-two subgroup classes in
$N\rtimes\langle s\rangle$.
Let $V_2=H^1(N,\Z_2)$, let
$e=q^{-1}\sum_{g\in R}g$ be the actual deck-averaging
operator, and put $W_2=(1-e)V_2$.
On this rank-$2(q-1)$ integral summand the positive
square acts identically:
\[
                         \beta^*|_{W_2}=I.                 \tag{D.69y}
\]
The statement determines the indicated quotient of the
centralizer of the positive square. It does not determine
the image of $C_F(s)$ or the full structure of either
centralizer, and it does not exclude or realize a negative
quartic coordinate.
\end{claim}
\begin{proof}
\emph{The actual covers and the two cyclic extensions.}
By D.68 the homomorphism $\rho$ has
\[
 \rho(x)=t,\qquad \rho(y)=r,\qquad
 \rho(c)=r^{-2},\qquad
 \rho\alpha=\rho,\qquad
 Q=\langle r,t\mid r^q=t^2=1,\ trt=r^{-1}\rangle .
                                                               \tag{D.69d}
\]
The last equality $\rho\alpha=\rho$ is the proved identity
action at this multiplier coordinate. Thus $N$ is normal
in $E=F\rtimes\langle s\rangle$, and the actual quotient,
with inverse coordinate maps, is
\[
 \begin{aligned}
 E/N&\longrightarrow Q\times C_4,
        &fs^jN&\longmapsto(\rho(f),s^j),\\
 Q\times C_4&\longrightarrow E/N,
        &(z,s^j)&\longmapsto fs^jN\quad(\rho(f)=z).
 \end{aligned}                                               \tag{D.69e}
\]
The second map is independent of the chosen lift $f$.
The multiplication rule follows from $\rho\alpha=\rho$,
so these maps are inverse homomorphisms.
Let
\[
 E_N=N\rtimes\langle s\rangle,\qquad
 E_{2,N}=N\rtimes\langle\sigma\rangle,\qquad
 F_0=\ker\chi=\rho^{-1}(R).
                                                               \tag{D.69f}
\]
These are the inverse images of the named factors under
(D.69e), with their actual generator labels.

The original peripheral monodromy $r^{-2}$ has order $q$.
Its right action on the $2q$ cosets has two cycles, both
of length $q$. Thus the finite covering surface has two
punctures. Its Euler characteristic is $-2q$, so its
genus is $q$. The inclusion of its discrete fundamental
group has finite index in the original free group and
therefore has its full induced profinite topology.
This identifies its full completion with $N$ and gives
$\operatorname{rank}(N)=2q+1$.
The intermediate group $F_0$ is the specified type-$(1,2)$
double cover with Schreier basis $u=x^2,B=y,C=xyx^{-1}$
from D.68. Its quotient $F_0/N=R$ sends these generators
to $1,r,r^{-1}$, respectively.

\emph{The exact integral deck module.}
Put $\Lambda=\Z_2[R]$ and $\nu=1+r+\cdots+r^{q-1}$.
The actual degree-$q$ graph cover of the three-petal rose
for $F_0$ has chain boundary
\[
 \partial:\Lambda^3\longrightarrow\Lambda,\qquad
       (a,b,d)\longmapsto(r-1)b+(r^{-1}-1)d .
                                                               \tag{D.69g}
\]
The two explicit inverse maps on its cycle module are
\[
 \begin{aligned}
 \Lambda^2\oplus\Z_2&\longrightarrow\ker\partial,
       &(a,d,k)&\longmapsto(a,r^{-1}d+k\nu,d),\\
 \ker\partial&\longrightarrow\Lambda^2\oplus\Z_2,
       &(a,b,d)&\longmapsto(a,d,k),
       \quad b-r^{-1}d=k\nu .
 \end{aligned}                                               \tag{D.69h}
\]
The scalar $k$ exists uniquely because the annihilator of
$r-1$ in $\Lambda$ consists exactly of constant-coefficient
multiples of $\nu$. This follows directly by comparing
the $q$ cyclic coefficients. Hence
\[
       H_1(N,\Z_2)\cong\Z_2\oplus\Lambda^2
                       \quad\text{as actual deck modules}. \tag{D.69i}
\]
Dualizing gives the same multiplicities in
$V_2=H^1(N,\Z_2)=\operatorname{Hom}_{\rm cont}(N,\Z_2)$.
Let $e=q^{-1}\sum_{g\in R}g$ and
\[
                V=H^1(N,\F_2),\qquad
                W_2=(1-e)V_2,\qquad W=W_2/2W_2 .
                                                               \tag{D.69j}
\]
Thus $eV_2$ has rank three and $W_2$ has rank $2(q-1)$.
Restriction and transfer for the index-$q$ inclusion
$N\subset F_0$, with $q$ invertible in $\Z_2$, identify
$eV_2$ with the actual $H^1(F_0,\Z_2)$. These maps are
equivariant for the original $\alpha$.

Let $S$ act on $V_2$ by pullback under $\alpha^{-1}$, and
write $B=S^2$ for the pullback under $\beta$.
Because $\alpha$ induces identity on $Q$, $S$ commutes
with every deck operator from $R$: the two actual
conjugations differ only by conjugation by an element
of $N$, which acts trivially on this homology.
Choose the finite unramified extension $\mathcal O$ of
$\Z_2$ containing the $q$th roots of unity. Its explicit
character projectors are
$q^{-1}\sum_{g\in R}\psi(g)^{-1}g$.
Every nontrivial character summand of
$W_2\otimes\mathcal O$ is free of rank two by (D.69i).
On its residue field $S^4=1$ makes $S$ unipotent.
A two-dimensional unipotent matrix in characteristic
two has square identity. Consequently
\[
                          B|_W=I.                          \tag{D.69k}
\]
On the deck-trivial rank-three summand, let
$S_{\rm hom}$ be the matrix induced by $\alpha$
on the retained homology basis of $F_0$.
D.68 gives $S_{\rm hom}^2=-J$.
In the dual basis the present pullback convention is
$S=(S_{\rm hom}^{-1})^{\mathsf T}$; therefore
$B=S^2=-J^{\mathsf T}$, since $J^2=I$.
The reduction of $J^{\mathsf T}-I$ has rank one.
Hence $S$ has one Jordan block of size three there,
and $B$ has blocks of sizes two and one.
On all of $V$, therefore,
\[
                   \operatorname{rank}(B-I)=1.             \tag{D.69l}
\]

\emph{Which full section classes occur in these covers.}
D.46 gives three $F$-conjugacy classes of order-two
sections over $\beta$, with zero mod-two first
cohomology for each of their full fixed groups.
The proof of D.53 gives their three distinct affine
labels in $H_1(F,\F_2)$. Use $\sigma=s^2$ as label zero.
Conjugation by $s$ now has zero translation part,
since it fixes $\sigma$ itself. Its linear part is
$A_2(a)$ modulo two, a nonidentity transvection.
Its unique fixed line is $\ker\chi$.
D.53 proves that its action fixes exactly the class
of $\sigma$ and interchanges the other two classes.
Their two labels consequently have $\chi$-value one.
An element $f\sigma$ in either class has
$\chi(f)=1$, so its image under (D.69e) has nontrivial
$Q$ coordinate. Neither class meets $E_{2,N}$.
Every finite subgroup of $E_N$ injects into its
quotient $C_4$, since $N$ is torsion-free.
Its order-two subgroups lie in $E_{2,N}$.
Thus the only global class meeting either group is
the class of $\langle\sigma\rangle$.

Put $C=C_F(\sigma)$. Its mod-two first cohomology
vanishes by D.46, so $\chi|_C=0$.
The image $\rho(C)$ is contained in $R$ and is
therefore either $1$ or $R$.
The exact class parametrization is
\[
 Q/\rho(C)\longrightarrow
 \{\text{$N$-conjugacy classes of sections in $E_{2,N}$}\},
 \qquad
 \rho(f)\rho(C)\longmapsto
          [\,f\langle\sigma\rangle f^{-1}\,].
                                                               \tag{D.69m}
\]
Indeed two such subgroups are $N$-conjugate precisely
when their conjugators differ by an element of $N$
and an element of $N_F(\langle\sigma\rangle)=C$.
Every conjugate belongs to $E_{2,N}$ because its
image in $Q\times C_4$ is the same central subgroup.
These are also its $E_{2,N}$-conjugacy classes:
multiply any conjugator by the matching element of
the section to leave a conjugator in $N$.
Conjugation by $s$ on (D.69m) sends $\rho(f)$ to
$\rho(\alpha(f))=\rho(f)$ and fixes $\sigma$.
Hence (D.69m) also parametrizes the $E_N$ classes.

\emph{The trivial-image alternative and its Brown algebra.}
Suppose $\rho(C)=1$.
Then $C\subset N$ and (D.69m) has exactly $2q$ classes.
The normalizer in $E_{2,N}$ of the representative
$f\langle\sigma\rangle f^{-1}$ is exactly
\[
                       f(C\times\langle\sigma\rangle)f^{-1}.
                                                               \tag{D.69n}
\]
The product map sends $(h,\sigma^\epsilon)$ to
$fh\sigma^\epsilon f^{-1}$; its inverse reads the
unique $F$ coordinate and the specified quotient
coordinate after conjugation by $f^{-1}$.
As a closed subgroup of the full free profinite
group $F$, $C$ has $\operatorname{cd}_2(C)\leq1$.
Since $H^1(C,\F_2)=0$, all its positive-degree
mod-two cohomology vanishes.

The full Brown restriction theorem
\cite[Theorem~2 and \S1.1]{A-KochloukovaZalesskii2007}
applies to $E_{2,N}$: its open subgroup $N$ has
$2$-cohomological dimension one, every finite subgroup
injects into $C_2$, and the only nontrivial elementary
abelian two-subgroups are the order-two sections,
with no inclusions between distinct ones.
The two-row Hochschild--Serre spectral sequence
has finite modules in every degree, since $N$
is finitely generated and $C_2$ is finite.
Thus $H^n(E_{2,N},\F_2)$ is finite in every degree;
the theorem's dense restriction embedding is
therefore onto its Hausdorff target.
The actual restriction map is a ring isomorphism
in sufficiently high degrees onto the product of
the cohomologies of (D.69n).
There are exactly $2q$ factors, each one-dimensional
in each positive degree.
For the quotient character $u:E_{2,N}\to\F_2$,
localizing that restriction and taking degree zero gives
\[
 \mathcal R=
 \bigl(H^*(E_{2,N},\F_2)[u^{-1}]\bigr)^0
                         \cong\F_2^{\,2q}.                \tag{D.69o}
\]
Here $u$ is zero on each fixed group and nonzero
on its order-two section. This is the actual
localization of the displayed restriction maps.

We need the following precise injection into this algebra.
Every $\theta\in W$ is invariant under $\beta$ by
(D.69k), and extends to a character of $E_{2,N}$
by declaring $\theta(\sigma)=0$.
The Hochschild--Serre spectral sequence has only
rows zero and one. The actual section
$\langle\sigma\rangle\to E_{2,N}$ makes inflation
from $H^*(C_2,\F_2)$ injective in every degree.
Thus every bottom-row class survives, and all
possible differentials, whose targets are in that
row, are zero. On $W$ the row-one module is
$W[u]$, since $\beta$ acts trivially modulo two.
The associated graded class of $u^m\theta$ is the
nonzero element $\theta u^m$ whenever $\theta\ne0$.
It follows that
\[
                W\longrightarrow\mathcal R,\qquad
                         \theta\longmapsto\theta/u
                    \quad\text{is injective}.             \tag{D.69p}
\]

\emph{The exact mod-four character cochain.}
Since $B=I$ on $W_2/2W_2$, the integral operator
$(I-B)/2$ is defined on $W_2$. Let $D$ be its
reduction modulo two.
Lift $\theta\in W$ to a continuous character
$\varphi:N\to\Z/4$ from the integral summand $W_2$.
The cochain
$\widetilde\theta(n\sigma^\epsilon)=\varphi(n)$
has exact coboundary
\[
 \begin{aligned}
 \delta\widetilde\theta
       (n\sigma^\epsilon,m\sigma^\eta)
 &=\varphi(n)+\varphi(m)
                   -\varphi(n\beta^\epsilon(m))\\
 &=\epsilon\bigl(\varphi(m)-\varphi(\beta(m))\bigr)
       =2\epsilon(D\theta)(m)\quad\text{in }\Z/4.
 \end{aligned}                                               \tag{D.69q}
\]
Changing the lift by twice a character does not
change $D\theta$, because $\beta$ is identity
modulo two on $W$.
Thus its Bockstein, which equals its cup square,
is the actual identity
\[
                           \theta^2=uD\theta .
                                                               \tag{D.69r}
\]
Every element of the product algebra (D.69o) is
idempotent. Applying this to $\theta/u$ and using
(D.69r) and the injection (D.69p) gives $D\theta=\theta$.
Consequently
\[
                         B|_{W_2}\equiv -I\pmod4.          \tag{D.69s}
\]

\emph{The incompatible quartic section count.}
Extend (D.69s) to the unramified ring $\mathcal O$
used above. On a nontrivial character summand,
if the two-dimensional reduction of $S$ were
identity, write $S=I+2L$ over $\mathcal O$.
Then $S^2\equiv I\pmod4$, contrary to
$S^2=B\equiv-I\pmod4$.
Hence every nontrivial character summand of
$V\otimes\overline{\F}_2$ has one Jordan block
of size two for $S$. The deck-trivial summand
has one block of size three, as already proved.
There are $q-1$ nontrivial characters.
In the periodic $C_4$ resolution the two maps
are $S-I$ and
$1+S+S^2+S^3=(S-I)^3$.
Each of these blocks contributes exactly one
dimension in every positive cohomological degree.
Descent of dimension from the splitting field gives
\[
                         \dim H^j(C_4,V)=q\quad(j>0).
                                                               \tag{D.69t}
\]
The two-row spectral sequence for $E_N$ therefore
gives the upper bound
\[
                         \dim H^n(E_N,\F_2)\leq q+1
                                      \quad(n\geq2).       \tag{D.69u}
\]

On the other hand, the exact full normalizer in
$E$ of $\langle\sigma\rangle$ is
$C\rtimes\langle s\rangle$, as in D.53.
For each representative in (D.69m), its conjugate
lies entirely in $E_N$ under the present
assumption $\rho(C)=1$: (D.69e) sends its fixed
group to identity and its conjugated section
generator to the named $C_4$ factor.
Thus its $E_N$ normalizer is exactly
\[
                           f(C\rtimes\langle s\rangle)f^{-1}.
                                                               \tag{D.69v}
\]
The map $(h,s^j)\mapsto fhs^jf^{-1}$ and its
coordinate inverse give this semidirect-product
identification with labels.
All positive-degree cohomology of $C$ is zero,
so inflation identifies the cohomology of each
normalizer with $H^*(C_4,\F_2)$.
Brown applies here as well: $N$ is the open
dimension-one subgroup, finite subgroups inject
into $C_4$, and there is no elementary abelian
subgroup of rank two. Its elementary two-subgroup
classes are exactly the $2q$ classes in (D.69m);
conjugation by $s$ does not fuse them.
The same two-row finite-module spectral sequence
makes $H^n(E_N,\F_2)$ finite, so this dense
restriction embedding is again onto.
Restriction therefore gives
$\dim H^n(E_N,\F_2)=2q$ in high degrees.
This contradicts (D.69u), since $q>1$.
The alternative $\rho(C)=1$ is impossible.
We have proved the surjectivity in (D.69b).

\emph{The exact remaining fixed-group cohomology.}
Now $\rho(C)=R$. Formula (D.69m) has exactly two
cosets, represented by $1$ and $t=\rho(x)$.
Intersecting the full normalizers with $E_{2,N}$
gives precisely (D.69c), since its fixed kernel is
$C_0=C\cap N$. Both are conjugate by the original
element $x$ and have the same first-cohomology dimension.

The split spectral sequence for $E_{2,N}$ still
collapses by the actual section argument above.
By (D.69l) the operator $B-I$ on its
$(2q+1)$-dimensional module $V$ has rank one.
For $C_2$ both periodic differentials are $B-I$,
so
\[
 \dim H^j(C_2,V)=(2q+1)-2=2q-1\quad(j>0),\qquad
 \dim H^n(E_{2,N},\F_2)=2q\quad(n\geq2).
                                                               \tag{D.69w}
\]
Brown restriction to the two normalizers (D.69c)
then gives the exact equality
\[
                    2q=2\bigl(1+\dim H^1(C_0,\F_2)\bigr).
                                                               \tag{D.69x}
\]
In particular this first cohomology is finite-dimensional,
and its dimension is $q-1$, as claimed.
No assertion that $C_0$ is itself a full free profinite
group, or that $C_F(s)$ has the same finite image,
is used.

\emph{The full integral action on the nontrivial deck summand.}
We now use the proved image $\rho(C)=R$, not the
trivial-image alternative discarded above.
Choose an actual $h\in C$ with $\rho(h)=r$.
Conjugation by $h$ preserves $E_{2,N}$ and fixes
$\sigma$ itself. Its $q$th power is conjugation
by $h^q\in C_0\subset E_{2,N}$, so the resulting
cohomology action factors through $R$.
On $V_2$ it is exactly the deck action: replacing
$h$ by any other lift of $r$ changes conjugation
by an inner automorphism of $N$.
The extension of $\theta\in W$ which is zero on
$\sigma$ is consequently $R$-equivariant, and
the quotient character $u$ is fixed.

By the same finite-dimensional Brown restriction
argument, the actual localized degree-zero ring is now
\[
 \mathcal R\cong
 \bigl(\F_2\oplus H^1(C_0,\F_2)u^{-1}\bigr)
 \times
 \bigl(\F_2\oplus H^1(xC_0x^{-1},\F_2)u^{-1}\bigr).
                                                               \tag{D.69z}
\]
In each factor the product of two degree-one fixed-group
classes vanishes, because its full fixed group has
$2$-cohomological dimension at most one.
Thus the two displayed augmentation ideals are square
zero; their product $\mathcal I$ is the nilradical of
$\mathcal R$, and $\mathcal R/\mathcal I=\F_2^2$.
The action of $R$ fixes both section classes, which
are the two cosets in $Q/R$. Its action on this
residue algebra is therefore identity.

The injection (D.69p) still holds, since its split
spectral-sequence proof used no assumption on
$\rho(C)$. It is now the $R$-equivariant map
$W\to\mathcal R$, $\theta\mapsto\theta/u$.
The composite with $\mathcal R/\mathcal I$ vanishes:
deck averaging is zero on $W$ and identity on that
residue algebra, and $q$ is invertible in $\F_2$.
Hence its image is contained in $\mathcal I$.
Both vector spaces have dimension $2(q-1)$ by
(D.69j) and (D.69x), so it identifies $W$ with
that square-zero ideal.
In particular $(\theta/u)^2=0$ for every $\theta\in W$.
The exact cochain identity (D.69r) and the same
injection give $D\theta=0$ for every $\theta$.
Thus $B|_{W_2}\equiv I\pmod4$.

Finally $B^2=I$ on the free $\Z_2$-module $W_2$.
Writing $B=I+4L$ gives
$L(I+2L)=0$ after cancelling the nonzero scalar
$8$ in its torsion-free endomorphism module.
The matrix $I+2L$ is invertible over $\Z_2$,
because its reduction is identity.
Therefore $L=0$ and $B=I$ on $W_2$.
This is precisely (D.69y), since $\beta^{-1}=\beta$.
Together with the preceding dual-basis calculation,
it determines the integral square as
$-J^{\mathsf T}$ on $eV_2$ and identity on $W_2$.

The executable certificate
\path{code/part_c/torus_quartic_bad_prime_d69.py}
checks the exact finite covering cycles and module maps,
the two section-label possibilities, the cyclic
resolution dimensions, and the mod-four cochain identity.
The full normalizer and Brown arguments above supply
the stated centralizer conclusion.
\end{proof}
